\centerline{\bf \Huge Рыцари и лжецы. Уравнения}
\title{ Рыцари и лжецы. Уравнения}

\problem{1} На острове живет 25 человек: рыцари, лжецы и хитрецы. Рыцари всегда говорят правду, лжецы всегда лгут, а хитрецы отвечают на заданные им вопросы по очереди то правду, то ложь. Все жителям острова было задано три вопроса: «Вы рыцарь?», «Вы хитрец?», «Вы лжец?». На первый вопрос «да» ответили 15 человек, на второй — 7 человек, на третий — 5 человек. Сколько хитрецов живет на этом острове?

\problem{2} Болельщики «Спартака» говорят правду, когда «Спартак»
выигрывает, и лгут, когда он проигрывает. Аналогично ведут себя болельщики «Динамо», «Зенита» и «Локомотива». После двух матчей с участием этих четырёх команд, каждая из которых закончилась победой одной из команд, а не ничьей, из болельщиков, смотревших трансляцию, на вопрос «Болеете ли вы за «Спартак»?» положительно ответили 200 человек, на вопрос «Болеете ли вы за «Динамо»?» положительно ответили 300 человек, на вопрос «Болеете ли вы за «Зенит»?» положительно ответили 500 человек, на вопрос «Болеете ли вы за «Локомотив»?» положительно ответили 600 человек. Сколько человек болело за каждую из команд?

\problem{3} В классе учатся 30 человек: отличники, троечники и двоечники. Отличники на все вопросы отвечают правильно, двоечники всегда ошибаются, а троечники на заданные им вопросы строго по очереди то отвечают верно, то ошибаются. Всем ученикам было задано по три вопроса: «Ты отличник?», «Ты троечник?», «Ты двоечник?». Ответили «Да»: на первый вопрос — 19 учащихся, на второй — 12, на третий — 9. Сколько троечников учится в этом классе?

\problem{4} В подземном царстве живут гномы, предпочитающие носить либо
зелёные, либо синие, либо красные кафтаны. Некоторые из них всегда лгут, а остальные всегда говорят правду. Однажды каждому из них задали четыре вопроса:\\
1. «Ты предпочитаешь носить зелёный кафтан?»\\
2. «Ты предпочитаешь носить синий кафтан?»\\
3. «Ты предпочитаешь носить красный кафтан?»\\
4. «На предыдущие вопросы ты отвечал честно?» \\
На первый вопрос «да» ответили 40 гномов, на второй — 50, на третий — 70, а на четвёртый — 100. Сколько честных гномов в подземном царстве?

\newpageУравнения

\problem{5} В некоторой школе каждый десятиклассник либо всегда говорит правду, либо всегда лжёт. Директор вызвал к себе нескольких десятиклассников и спросил каждого из них про каждого из остальных, правдивец тот или лжец. Всего было получено 44 ответа «правдивец» и 28 ответов «лжец». Сколько правдивых ответов мог получить директор?

\problem{6} В некотором классе каждый ученик либо всегда говорит ложь, либо всегда говорит правду. При этом каждый из них знает про остальных, кто лжец, а кто — нет. На сегодняшнем собрании присутствовали все ученики класса, и каждый сообщил, кем является каждый из остальных. Ответ «лжец» при этом прозвучал 272 раза. Вчера проводилось такое же собрание, но один из учеников отсутствовал. Тогда ответ «лжец» прозвучал 256 раз. Сколько учеников в таком классе?
  
\problem{7} В городе 200 жителей. Часть из них — рыцари, которые всегда говорят
правду, остальные — лжецы, которые всегда лгут. Каждый горожанин живет в одном из четырёх кварталов (А, Б, В и Г). Каждому задали четыре вопроса: «Вы живёте в квартале А?», «Вы
живёте в квартале Б?», «Вы живёте в квартале В?», «Вы живёте в квартале Г?». На первый вопрос утвердительно ответили 105 жителей, на второй — 45, на третий — 85 и на четвёртый — 65.
В каком квартале лжецов живет больше, чем рыцарей, и на сколько?